\begin{table}[t]
\centering
\small
\begin{tabular}{llccc}
\toprule
 & & $\alpha = 0.1$ & $\alpha = 0.5$ & $\alpha = 100$ \\
\midrule
FedRR & Healthy runs with any alarm & 6 of 6 & 6 of 6 & 3 of 6 \\
 & False alarms, \% per client per round & 13.9 & 14.1 & 2.4 \\
 & Fault runs caught & 18 of 60 & 57 of 84 & 78 of 84 \\
\addlinespace
step-normalised & Healthy runs with any alarm & 6 of 6 & 6 of 6 & 6 of 6 \\
 & False alarms, \% per client per round & 13.3 & 14.4 & 8.8 \\
 & Fault runs caught & 60 of 60 & 75 of 84 & 78 of 84 \\
\bottomrule
\end{tabular}
\caption{FedRR at its own control limit ($H = 3.86$ for $K = 10$, $d = 0.5$, in-control average run length 30, as in its paper). No run is used for calibration, so all six healthy runs per level count. Step-normalised: the residual divided by the client's number of local SGD steps that round, a rough correction for clients with more data taking more steps; it does not make clients exchangeable. Fault runs: every fault run FedRR could be scored on at that level, all severities.}
\label{tab:fedrr}
\end{table}
